\chapter{装进实体机}
\label{ch:physical-install}
\chapterlead{实体机安装多了三件事：确认磁盘名、处理引导与固件、为笔记本补齐无线与休眠。本章在虚拟机流程之上补充这些差异，并说明双系统与磁盘加密的取舍。}

\section{起步：制作安装介质}

需要一个 4\,GB 以上的 U 盘，内容会被清空。

\begin{shellsession}
$ lsblk                      # 确认 U 盘设备名，例如 /dev/sdb
$ sudo dd if=nixos-gnome-26.11-x86_64-linux.iso of=/dev/sdb bs=4M status=progress oflag=sync
\end{shellsession}

\texttt{of=} 必须写整块设备（\texttt{/dev/sdb}），不是分区（\texttt{/dev/sdb1}）。写错目标会毁掉那块磁盘上的数据。Windows 上可用 Rufus 或 balenaEtcher，macOS 上同样用 \texttt{dd}，设备名为 \texttt{/dev/rdisk2} 一类。

\section{确认磁盘与固件}

\begin{shellsession}
$ lsblk -o NAME,SIZE,MODEL,SERIAL
$ ls /sys/firmware/efi && echo "UEFI 启动"
$ mokutil --sb-state
\end{shellsession}

\texttt{MODEL} 与 \texttt{SERIAL} 用于确认目标盘，避免误操作。\texttt{/sys/firmware/efi} 存在说明当前以 UEFI 模式启动，这是 NixOS 默认引导方式的前提。\texttt{mokutil} 报告 Secure Boot 状态；开启时需要额外处理签名，初次安装建议先在固件设置里关闭。

\begin{warnbox}[先备份再分区]
实体机上的分区操作不可撤销。安装前把原有数据复制到外部磁盘，并确认备份可以读取。安装了双系统时，尤其要确认 Windows 的快速启动与休眠已关闭，否则它的分区会处于不一致状态。
\end{warnbox}

\section{分区方案}

虚拟机章节用了最简单的两分区方案。实体机上按需要选择。

\begin{booktable}{\textwidth}{J{0.846} J{1.154}}
\toprule
方案 & 组成 \\
\midrule
\endhead
最简 & EFI 512\,MB + 根分区（余下全部）\\
含交换 & EFI 512\,MB + 交换（内存大小）+ 根分区 \\
分离 home & EFI 512\,MB + 根 40--60\,GB + home（余下）\\
加密 & EFI 512\,MB + LUKS 容器，容器内含交换与根 \\
\bottomrule
\end{booktable}

分离 \texttt{/home} 的好处是重装系统时保留用户数据；代价是分区大小需要预先规划。加密适合笔记本：设备丢失时数据不可读，代价是每次开机要输入口令，且无法远程重启后自动解锁（可配合 SSH 解锁方案，但增加复杂度）。

加密安装的骨架：

\begin{shellsession}
$ sudo cryptsetup luksFormat /dev/nvme0n1p2
$ sudo cryptsetup open /dev/nvme0n1p2 cryptroot
$ sudo mkfs.ext4 /dev/mapper/cryptroot
$ sudo mount /dev/mapper/cryptroot /mnt
\end{shellsession}

生成的 \texttt{hardware-configuration.nix} 会包含解密配置；若没有，需要手工补 \texttt{boot.initrd.luks.devices}。

\section{双系统}

与 Windows 共存时，保留它的 EFI 分区，把 NixOS 装在空闲分区上。

\begin{nixcode}
boot.loader = {
  systemd-boot.enable = true;
  efi.canTouchEfiVariables = true;
  systemd-boot.windows = {
    "windows" = {
      title = "Windows";
      efiDeviceHandle = "HD0b";
    };
  };
};
\end{nixcode}

\texttt{efiDeviceHandle} 需要按实际情况填写，可用 \texttt{bootctl list} 查看。另一种做法是在固件启动菜单里切换，不改配置。

Windows 更新有时会重写引导顺序，导致直接进入 Windows。此时进固件设置把 NixOS 调回第一位，或在 Windows 中用 \texttt{bcdedit} 调整。

\section{笔记本的额外步骤}

装完后补齐无线、蓝牙、电源与背光。

\begin{nixcode}
hardware.enableRedistributableFirmware = true;
hardware.bluetooth.enable = true;
services.tlp.enable = true;
services.fwupd.enable = true;

# 部分机型需要显式声明内核模块
boot.kernelModules = [ "kvm-intel" ];
boot.kernelParams = [ "mem_sleep_default=deep" ];
\end{nixcode}

\texttt{fwupd} 用于固件更新，执行 \texttt{fwupdmgr get-updates} 查看可更新的固件。\texttt{mem\_sleep\_default=deep} 让部分机型进入真正的休眠而不是待机耗电状态。

\begin{shellsession}
$ sudo nixos-generate-config --show-hardware-config > /tmp/hw.nix
$ diff /tmp/hw.nix /etc/nixos/hardware-configuration.nix
\end{shellsession}

重装或更换硬件后重新生成硬件配置，用 \texttt{diff} 对比可以看清哪些设备变了。

\section{练习}

\begin{enumerate}
  \item 用 \texttt{lsblk -o NAME,SIZE,MODEL,SERIAL} 确认目标磁盘，记录型号与序列号。
  \item 用 \texttt{dd} 制作安装介质，写完后用 \texttt{sha256sum} 校验 U 盘内容与镜像一致。
  \item 检查本机是否以 UEFI 启动，Secure Boot 处于什么状态。
  \item 按“分离 home”方案规划分区，写出每一步命令。
  \item 装完后用 \texttt{fwupdmgr get-devices} 列出可更新固件的设备。
\end{enumerate}

\begin{checkpoint}
本章结束时，应能：
\begin{itemize}
  \item 用 \texttt{dd} 制作安装介质，知道写错目标的后果；
  \item 确认磁盘型号与固件模式，判断能否使用默认引导方案；
  \item 在四种分区方案中选择一种，说明各自代价；
  \item 与 Windows 共存并处理引导顺序被改写的情况；
  \item 为笔记本补齐无线、蓝牙、电源与固件更新。
\end{itemize}
\end{checkpoint}
